\section{The Signal Was Real. The Strategy Was Not.}

The central finding is the distinction between \textit{what can be predicted} and \textit{what can be traded}.

\begin{table}[htbp]
\centering
\caption{Decision gates and outcomes.}
\label{tab:gates}
\begin{tabular}{lcc}
\toprule
\textbf{Criterion} & \textbf{Result} & \textbf{Gate} \\
\midrule
Predictive ranking & Passed & IC $> 0$, NW $t > 2$ \\
Time-aware significance & Passed & Bootstrap CI excludes 0 \\
Within-Q1 selection & Passed & Excess $> 0$ \\
Positive U50 return & Failed & Net $< 0$ \\
Acceptable drawdown & Failed & DD $> 25\%$ \\
Scalable capacity & Failed & Q1 holds +48.1 \\
\midrule
\textbf{Decision} & \textbf{TERMINATED} & OHLCV branch closed \\
\bottomrule
\end{tabular}
\end{table}

The frozen baseline retained a mean out-of-sample IC of 0.108 with Newey--West $t = 10.5$ and a block-bootstrap 95\% confidence interval lower bound of 0.087.

\subsection{Prediction vs.\ Decision}
Cross-sectional return prediction---even statistically robust prediction---is not the same as portfolio construction. In this project, each step revealed a new constraint: statistical significance survived all corrections; execution introduced timing gaps; costs consumed 4+ percentage points of annual return; liquidity concentrated the signal in stocks too small to trade at scale; and drawdown remained above threshold.

\subsection{The Discipline of Stopping}
The pre-registered decision gates were designed specifically to prevent post-hoc optimization. When all rescue attempts failed the same gates, the only honest decision was termination. This is not a failure of the research program---it is the program's primary output.
